
\begin{filecontents}{refs.bib}
    @article{Achyuth2026What,
    title={What Can Low Resource Languages Learn From Each Other?},
    author={P. Achyuth and Kahaan Shah and Chetan Arora},
    year={2026},
    doi={10.48550/arxiv.2608.27753}
    }
    
    @article{Al-Azzawi2026Cross-Lingual,
    title={Cross-Lingual Learning within Arabic Script for Low-Resource HTR},
    author={S. Al-Azzawi and E. Barney and M. Liwicki},
    year={2026},
    doi={10.48550/arxiv.2605.02089}
    }
    
    @article{Al-Azzawi2026Understanding,
    title={Understanding Cross-Language Transfer Improvements in Low-Resource HTR: The Role of Sequence Modeling},
    author={S. Al-Azzawi and Chang Liu and Nudrat Habib and E. Barney and Marcus Liwicki},
    journal={ArXiv},
    year={2026},
    volume={abs/2605.05900},
    doi={10.48550/arxiv.2605.05900}
    }
    
    @article{Al-Azzawi2026A,
    title={A human-in-the-loop label error detection framework applied to Arabic-script HTR datasets},
    author={S. Al-Azzawi and E. Barney and György Kovács and Marcus Liwicki},
    journal={Pattern Recognition},
    year={2026},
    doi={10.1016/j.patcog.2026.114621}
    }
    
    @article{Al-Azzawi2026Performance,
    title={Performance Gap Analysis between Latin and Arabic Scripts HTR},
    author={S. Al-Azzawi and E. Barney and M. Liwicki},
    journal={ArXiv},
    year={2026},
    volume={abs/2606.18884},
    doi={10.48550/arxiv.2606.18884}
    }
    
    @article{Alhomoud2025Cross-Lingual,
    title={Cross-Lingual SynthDocs: A Large-Scale Synthetic Corpus for Any to Arabic OCR and Document Understanding},
    author={Haneen A. Alhomoud and Asma A. Ibrahim and Murtadha Al-Jubran and Fahad Al-Otaibi and Y. Alharbi and D. Toibazar and Ke-Sen Wang and Pedro J. Moreno},
    journal={2025 IEEE International Conference on Data and Software Engineering (ICoDSE)},
    year={2025},
    pages={142-147},
    doi={10.1109/icodse68111.2025.11351903}
    }
    
    @article{Bourne2025Scrambled,
    title={Scrambled text: fine-tuning language models for OCR error correction using synthetic data},
    author={Jonathan Bourne},
    journal={International Journal on Document Analysis and Recognition (IJDAR)},
    year={2025},
    volume={28},
    pages={741 - 755},
    doi={10.1007/s10032-025-00522-0}
    }
    
    @article{Chandrashekar2025Applying,
    title={Applying Few-Shot Learning to Optical Character Recognition in Low-Resource Languages},
    author={B. Chandrashekar and Yogita D. Bhise and D. Muthukrishnaveni and Arudra Annepu and Swapnil M. Parikh and K. C.},
    journal={2025 2nd International Conference on Intelligent Algorithms for Computational Intelligence Systems (IACIS)},
    year={2025},
    pages={1-6},
    doi={10.1109/iacis65746.2025.11211315}
    }
    
    @article{Chung2025Finetuning,
    title={Finetuning Vision-Language Models as OCR Systems for Low-Resource Languages: A Case Study of Manchu},
    author={Yan Hon Michael Chung and Dong-Hun Choi},
    journal={ArXiv},
    year={2025},
    volume={abs/2507.06761},
    doi={10.48550/arxiv.2507.06761}
    }
    
    @article{Chung2026Fine-tuning,
    title={Fine-tuning Vision-Language Models as OCR Systems for Low-Resource Languages: A Case Study of Manchu},
    author={Yan Hon Michael Chung and Dong-Hun Choi},
    journal={Computational Humanities Research},
    year={2026},
    doi={10.1017/chr.2026.10042}
    }
    
    @article{Elkousy2026Domain-Specific,
    title={Domain-Specific Adaptation of Vision-Language Models for Arabic OCR},
    author={Amr Elkousy and Youssef Elasrigy and A. Ammar and M. Ibrahim and Mariam M. N. Aboelwafa},
    year={2026},
    pages={351-358},
    doi={10.1016/j.procs.2026.01.043}
    }
    
    @article{Evani2025SemiHastakshar:,
    title={SemiHastakshar: Generalizable Indic Handwritten OCR through Semi-Supervised Learning},
    author={Lalitha Evani and Ajoy Mondal and Jawahar C. V.},
    journal={Proceedings of the Sixteen Indian Conference on Computer Vision, Graphics and Image Processing},
    year={2025},
    doi={10.1145/3774521.3774605}
    }
    
    @article{Faraz2026Designing,
    title={Designing Production-Scale OCR for India: Multilingual and Domain-Specific Systems},
    author={Ali Faraz and Raja Kolla and Ashish Kulkarni and S. Agarwal},
    journal={ArXiv},
    year={2026},
    volume={abs/2602.16430},
    doi={10.48550/arxiv.2602.16430}
    }
    
    @article{Farazi2026When,
    title={When Do VLMs Help Arabic Manuscript OCR? A Cross-Dataset Study},
    author={M. Farazi and Firoj Alam and A. Maaradji and Zakaria Maamar and Hamdy Mubarak and Wajdi Zaghouani},
    year={2026},
    doi={10.48550/arxiv.2608.22366}
    }
    
    @article{Guechaoui2026RefLAM:,
    title={RefLAM: A Reference-Grounded Line Annotation Pipeline for Historical Arabic Manuscripts},
    author={Mohamed Guechaoui and Mohamed Diaa Zellagui and Souleyman Chaib and S. Dhelim},
    year={2026},
    doi={10.48550/arxiv.2608.25140}
    }
    
    @article{Gunda2026CURIO:,
    title={CURIO: Curvature-Aligned and Efficient OCR for Low-Resource Historical Manuscripts},
    author={Sai Madhusudan Gunda and Tathagata Ghosh and Simran Singh Sandral and Ravi Kiran Sarvadevabhatla},
    journal={2026 IEEE/CVF Winter Conference on Applications of Computer Vision (WACV)},
    year={2026},
    pages={2011-2021},
    doi={10.1109/wacv61042.2026.00200}
    }
    
    @article{Haq2026PsOCR:,
    title={PsOCR: Benchmarking large multimodal models for optical character recognition in low-resource pashto language},
    author={Ijazul Haq and Yingjie Zhang and Muhammad Saqib},
    journal={Ain Shams Engineering Journal},
    year={2026},
    doi={10.1016/j.asej.2026.104024}
    }
    
    @article{Heakl2026DocAtlas:,
    title={DocAtlas: Multilingual Document Understanding Across 80+ Languages},
    author={Ahmed Heakl and Youssef Mohamed and M. Sohail and R. Elbadry and A. Nassar and P. Staar and F. Khan and Imran Razzak and Salman H. Khan},
    journal={ArXiv},
    year={2026},
    volume={abs/2605.12623},
    doi={10.48550/arxiv.2605.12623}
    }
    
    @article{Heakl2025KITAB-Bench:,
    title={KITAB-Bench: A Comprehensive Multi-Domain Benchmark for Arabic OCR and Document Understanding},
    author={Ahmed Heakl and M. Sohail and Mukul Ranjan and Rania Hossam and Ghazi Ahmed and Mohamed El-Geish and Omar Maher and Zhiqiang Shen and F. Khan and Salman H. Khan},
    journal={ArXiv},
    year={2025},
    volume={abs/2502.14949},
    doi={10.48550/arxiv.2502.14949}
    }
    
    @article{Hennara2025Baseer:,
    title={Baseer: A Vision-Language Model for Arabic Document-to-Markdown OCR},
    author={Khalil Hennara and Muhammad Hreden and Mohamed Motasim Hamed and Ahmad Bastati and Zeina Aldallal and Sara Chrouf and Safwan AlModhayan},
    journal={ArXiv},
    year={2025},
    volume={abs/2509.18174},
    doi={10.48550/arxiv.2509.18174}
    }
    
    @article{Kargaran2026GlotOCR,
    title={GlotOCR Bench: OCR Models Still Struggle Beyond a Handful of Unicode Scripts},
    author={Amir Hossein Kargaran and Nafiseh Nikeghbal and Jana Diesner and Franccois Yvon and Hinrich Schutze},
    journal={ArXiv},
    year={2026},
    volume={abs/2604.12978},
    doi={10.48550/arxiv.2604.12978}
    }
    
    @article{Li2026HunyuanOCR-1.5:,
    title={HunyuanOCR-1.5: Making Lightweight OCR VLMs Faster and Better},
    author={Gengluo Li and Xingyu Wan and Shangpin Peng and Wei-Nong Wang and Hao Feng and Yongkun Du and Bing-Hong Wu and Zheng-Hao Ruan and Zhiqiong Lu and Liang Wu and Pengyuan Lyu and Huawen Shen and Zibin Lin and Shijing Hu and Jie Yang and Hongbing Wen and Guang-Hua Yu and Hong Liu and Bochao Wang and Can Ma and Han Hu and Chengquan Zhang and Yu ZHOU},
    journal={ArXiv},
    year={2026},
    volume={abs/2607.04884},
    doi={10.48550/arxiv.2607.04884}
    }
    
    @article{Mahdi2026Persian,
    title={Persian Pixel: A large-scale synthetic OCR dataset for Persian language},
    author={Pouria Mahdi and Haq Nawaz Malik},
    journal={ArXiv},
    year={2026},
    volume={abs/2607.20385},
    doi={10.48550/arxiv.2607.20385}
    }
    
    @article{Malik2026synthocr-gen:,
    title={synthocr-gen: A synthetic ocr dataset generator for low-resource languages- breaking the data barrier},
    author={Haq Nawaz Malik and K. Shafi and Tanveer Ahmad Reshi},
    journal={ArXiv},
    year={2026},
    volume={abs/2601.16113},
    doi={10.48550/arxiv.2601.16113}
    }
    
    @article{Malik2026Koshur,
    title={Koshur Pixel: a large-scale synthetic ocr dataset for kashmiri},
    author={Haq Nawaz Malik and Faiza Iqbal and Nahfid Nissar},
    journal={ArXiv},
    year={2026},
    volume={abs/2606.23144},
    doi={10.48550/arxiv.2606.23144}
    }
    
    @article{Malik2026600k-ks-ocr:,
    title={600k-ks-ocr: a Large-scale Synthetic Dataset for Optical Character Recognition in Kashmiri Script},
    author={Haq Nawaz Malik},
    journal={ArXiv},
    year={2026},
    volume={abs/2601.01088},
    doi={10.48550/arxiv.2601.01088}
    }
    
    @article{Manna2025Label-Free,
    title={Label-Free Adaptation of Indic Printed OCR},
    author={Akash Manna and Radha Krishna Deshpande and Ajoy Mondal and C. V. Jawahar},
    journal={Proceedings of the Sixteen Indian Conference on Computer Vision, Graphics and Image Processing},
    year={2025},
    doi={10.1145/3774521.3774533}
    }
    
    @article{Mathew2022Towards,
    title={Towards Deployable OCR Models for Indic Languages},
    author={Minesh Mathew and C.V. Jawahar},
    year={2022},
    pages={167-182},
    doi={10.1007/978-3-031-78495-8_11}
    }
    
    @article{Momayiz2026AtlasOCR:,
    title={AtlasOCR: Building the First Open-Source Darija OCR Model with Vision Language Models},
    author={Imane Momayiz and Soufiane Ait Elaouad and Abdeljalil Elmajjodi and Haitame Bouanane},
    journal={ArXiv},
    year={2026},
    volume={abs/2604.08070},
    doi={10.48550/arxiv.2604.08070}
    }
    
    @article{Saini2022OCR,
    title={OCR Synthetic Benchmark Dataset for Indic Languages},
    author={Naresh Saini and Promodh Pinto and Aravinth Bheemaraj and Deepak Kumar and Dhiraj Daga and Saurabh Yadav and Srihari Nagaraj},
    journal={ArXiv},
    year={2022},
    volume={abs/2205.02543},
    doi={10.48550/arxiv.2205.02543}
    }
    
    @article{Singh2026Can,
    title={Can OCR-VLMs Read Devanagari? A Stress-Test Benchmark and Post-Correction Study},
    author={A. Singh},
    journal={ArXiv},
    year={2026},
    volume={abs/2606.29213},
    doi={10.48550/arxiv.2606.29213}
    }
    
    @article{Taghadouini2026LightOnOCR:,
    title={LightOnOCR: A 1B End-to-End Multilingual Vision-Language Model for State-of-the-Art OCR},
    author={Said Taghadouini and A. Cavaillès and Baptiste Aubertin},
    journal={ArXiv},
    year={2026},
    volume={abs/2601.14251},
    doi={10.48550/arxiv.2601.14251}
    }
    
    @article{Wasfy2025QARI-OCR:,
    title={QARI-OCR: High-Fidelity Arabic Text Recognition through Multimodal Large Language Model Adaptation},
    author={Ahmed Wasfy and Omer Nacar and Abdelakreem Elkhateb and Mahmoud Reda and Omar Elshehy and A. Ammar and Wadii Boulila},
    journal={ArXiv},
    year={2025},
    volume={abs/2506.02295},
    doi={10.48550/arxiv.2506.02295}
    }
    
    @article{Xu2026Multilingual,
    title={Multilingual OCR-Aware Fine-Tuning and Prompt-Guided Chain-of-Thought Reasoning for Multimodal Large Language Models},
    author={Qin-Wu Xu and Yi-Fan Jiang and Haoyu Ren},
    journal={ArXiv},
    year={2026},
    volume={abs/2605.16409},
    doi={10.48550/arxiv.2605.16409}
    }
    \end{filecontents}
    
    \documentclass{article}
    \usepackage{adjustbox}
    \begin{document}
    
    \section*{\textbf{Small OCR VLM} fine-tuning lowers CER, but \textbf{real-scan robustness} and \textbf{script coverage} still dominate outcomes.}
    
    
    \subsection*{1. Introduction}
    
    
    Low-resource Indic and Arabic-script OCR has improved sharply since 2024, but the gains are uneven across training method, script, and evaluation setting. Small-to-mid-scale vision-language OCR models benefit from parameter-efficient adaptation, synthetic pretraining, and cross-script transfer, yet they still fail abruptly when evaluated on real scanned pages, degraded layouts, or scripts underrepresented in pretraining  \cite{Elkousy2026Domain-Specific} \cite{Wasfy2025QARI-OCR:} \cite{Faraz2026Designing} \cite{Singh2026Can} \cite{Kargaran2026GlotOCR}. Arabic-specific adaptation has been particularly successful: QARI-OCR, built from Qwen2-VL-2B with LoRA and optional 4-bit quantization, reached 6.1\% CER on diacritically rich Arabic text, while domain-adapted Qwen2.5-VL models cut CER on modern Arabic print by 29\%  \cite{Wasfy2025QARI-OCR:} \cite{Elkousy2026Domain-Specific}. Indic results are more mixed: specialized multilingual systems now reach near-state-of-the-art across several scripts, but clean synthetic text often exaggerates performance and real printed or handwritten scans remain much harder  \cite{Faraz2026Designing} \cite{Achyuth2026What} \cite{Singh2026Can} \cite{Mathew2022Towards}.
    
    
    Across the corpus, the main determinants of CER are not any single trick, but the interaction of four factors: whether the base model already has OCR priors, whether training data mixes realistic synthetic and real scans, whether adaptation preserves multilingual coverage without forgetting, and whether evaluation uses authentic degraded documents rather than rendered pages  \cite{Faraz2026Designing} \cite{Hennara2025Baseer:} \cite{Heakl2026DocAtlas:} \cite{Alhomoud2025Cross-Lingual} \cite{Singh2026Can}. Evidence for hard-example mining and reinforcement learning with verifiable rewards is promising but still narrow, with stronger support for reducing failure modes and hallucinations than for proven CER gains on low-resource Indic or Arabic scripts specifically  \cite{Taghadouini2026LightOnOCR:} \cite{Xu2026Multilingual}.
    
    
    \textbf{Figure 1:} Consensus on efficient adaptation for low-resource OCR
    
    
    \subsection*{2. Methods}
    
    
    This review synthesizes the supplied Deep Search corpus on low-resource OCR adaptation for Indic and Arabic-script languages, emphasizing evidence on LoRA or QLoRA fine-tuning, synthetic versus real training data, quantization, reinforcement learning, labelled-data efficiency, and comparisons with 2024–2026 multilingual OCR systems.
    
    
    The search covered the Consensus literature index spanning more than 220 million research papers from Semantic Scholar, PubMed, and other sources, then screened for papers directly reporting OCR or HTR outcomes, adaptation strategies, and benchmark comparisons relevant to low-resource scripts.
    
    
    \textbf{Figure 2:} Literature search and screening funnel for OCR studies
    
    
    The strategy combined targeted searches on OCR adaptation, synthetic-to-real transfer, quantization, RL-based optimization, and multilingual low-resource benchmarking.
    
    
    \subsection*{3. Results}
    
    
    \subsubsection*{3.1 Key Papers}
    
    
    A small set of papers anchors the field because they directly test the methods asked about under realistic multilingual OCR settings. The most influential recent anchors are QARI-OCR for Arabic LoRA adaptation, Chitrapathak for Indic multilingual production OCR, DocAtlas for multilingual adaptation trade-offs, KITAB-Bench and GlotOCR Bench for cross-system evaluation, and LightOnOCR for compact end-to-end OCR VLM design  \cite{Wasfy2025QARI-OCR:} \cite{Faraz2026Designing} \cite{Heakl2026DocAtlas:} \cite{Heakl2025KITAB-Bench:} \cite{Kargaran2026GlotOCR} \cite{Taghadouini2026LightOnOCR:}.
    
    
    \mbox{}\\
    \begin{adjustbox}{max width=\textwidth}
    \begin{tabular}{p{8.0cm} | p{8.0cm}}
    \hline
    Paper & Summary \\
    \hline
    \cite{Wasfy2025QARI-OCR:} & Arabic 2B VLM reaches 6.1\% CER  \cite{Wasfy2025QARI-OCR:} \\
    \cite{Faraz2026Designing} & OCR-specialized fine-tuning beats generic VLM training  \cite{Faraz2026Designing} \\
    \cite{Heakl2026DocAtlas:} & DPO and QKV-LoRA best preserve multilingual transfer  \cite{Heakl2026DocAtlas:} \\
    \hline
    \end{tabular}
    \end{adjustbox}
    
    
    \textbf{Figure 3:} Key papers anchoring recent low-resource OCR literature
    
    
    \subsubsection*{3.2 Adaptation, Quantization, and Data Efficiency}
    
    
    LoRA or QLoRA fine-tuning consistently improves low-resource OCR when the base model already has strong OCR or multilingual visual-text priors  \cite{Elkousy2026Domain-Specific} \cite{Momayiz2026AtlasOCR:} \cite{Faraz2026Designing} \cite{Xu2026Multilingual}. In Arabic, 4-bit LoRA on Qwen2.5-VL reduced CER by 29\% on modern print and improved historical-document accuracy, outperforming standard OCR engines on KITAB-Bench  \cite{Elkousy2026Domain-Specific}. QARI-OCR inserted rank-16 LoRA adapters into both vision and language modules, optionally with 4-bit quantization, and still achieved open-source state-of-the-art Arabic CER  \cite{Wasfy2025QARI-OCR:}. AtlasOCR extended this pattern to a 3B Qwen2.5-VL model with QLoRA, showing that a compact Arabic-family OCR model can challenge larger systems after targeted synthetic-plus-real adaptation  \cite{Momayiz2026AtlasOCR:}.
    
    
    Data efficiency depends strongly on initialization. In Indic printed OCR, same-language pretraining reached 92\% WRR, cross-lingual pretraining 51\%, and random initialization 6\%, showing that adaptation quality matters more than merely adding trainable parameters  \cite{Manna2025Label-Free}. Chitrapathak-2, which fine-tunes an OCR-specialized model, outperformed an end-to-end generic VLM trained on more data, indicating that OCR priors materially lower labelled-data needs  \cite{Faraz2026Designing}. Cross-script training also helps under low-resource regimes: Arabic-script HTR improved at 100 to 1000 labelled lines, and Persian reached 9.99 CER without full-scale target-language data  \cite{Al-Azzawi2026Cross-Lingual}.
    
    
    \mbox{}\\
    \begin{adjustbox}{max width=\textwidth}
    \begin{tabular}{p{4.0cm} | p{4.0cm} | p{4.0cm} | p{4.0cm}}
    \hline
    Dimension & Stronger Outcome & Weaker Outcome & Citations \\
    \hline
    Initialization & OCR-specialized or same-script & Random scratch training & \cite{Faraz2026Designing} \cite{Manna2025Label-Free} \\
    Quantized PEFT & 4-bit LoRA often preserves gains & Full effect vs full FT still task-dependent & \cite{Wasfy2025QARI-OCR:} \cite{Elkousy2026Domain-Specific} \cite{Faraz2026Designing} \\
    Very low data & Cross-lingual joint training helps & Benefits shrink as target data grows & \cite{Al-Azzawi2026Cross-Lingual} \cite{Al-Azzawi2026Understanding} \\
    Label needs & 100–1000 lines can help transfer & <10k real images often still limiting & \cite{Al-Azzawi2026Cross-Lingual} \cite{Achyuth2026What} \\
    \hline
    \end{tabular}
    \end{adjustbox}
    
    
    \textbf{Figure 4:} Adaptation trade-offs for low-resource OCR fine-tuning
    
    
    Key difference: 4-bit LoRA is now credible for OCR adaptation, but the largest gains come from model choice and training data realism rather than quantization alone  \cite{Elkousy2026Domain-Specific} \cite{Bourne2025Scrambled} \cite{Hennara2025Baseer:}.
    
    
    \subsubsection*{3.3 Synthetic Versus Real Pages}
    
    
    Synthetic data is essential for low-resource scripts because manual annotation remains the main bottleneck  \cite{Malik2026synthocr-gen:} \cite{Saini2022OCR} \cite{Mahdi2026Persian} \cite{Guechaoui2026RefLAM:}. Large synthetic corpora now exist for Pashto, Kashmiri, Persian, and multi-script Indic OCR, often with hundreds of thousands to millions of samples and degradation pipelines intended to narrow the synthetic–real gap  \cite{Haq2026PsOCR:} \cite{Malik2026synthocr-gen:} \cite{Malik2026Koshur} \cite{Mahdi2026Persian} \cite{Malik2026600k-ks-ocr:}. Synthetic-only training can transfer surprisingly well when degradations are realistic: Manchu VLMs trained on 60,000 synthetic words retained 87\% to 93\% real-document word accuracy, far better than a CRNN baseline that collapsed on scans  \cite{Chung2025Finetuning} \cite{Chung2026Fine-tuning}.
    
    
    But synthetic benchmarks often overstate OCR quality. On Devanagari, all ten systems clustered tightly on clean renders, yet nine of ten dropped sharply on real scans, with EasyOCR falling from chrF++ 93.6 to 58.3 and a 76-point spread opening across systems  \cite{Singh2026Can}. Arabic studies report the same pattern: synthetic corpora help, especially when scanned backgrounds and diacritic-aware fonts are included, but mixed synthetic-plus-real training is more robust than synthetic alone  \cite{Alhomoud2025Cross-Lingual} \cite{Hennara2025Baseer:}. Baseer’s 300k synthetic plus 200k real mixture and Cross-Lingual SynthDocs both outperformed base Qwen variants and narrowed the gap to Gemini on Arabic OCR benchmarks  \cite{Hennara2025Baseer:} \cite{Alhomoud2025Cross-Lingual}.
    
    
    \begin{itemize}
    \item \textbf{May 2022}\begin{itemize}
    \item 2 papers:  \cite{Mathew2022Towards} \cite{Saini2022OCR}- \textbf{Feb 2025}
    \item 1 paper:  \cite{Heakl2025KITAB-Bench:}- \textbf{Apr 2025}
    \item 1 paper:  \cite{Bourne2025Scrambled}- \textbf{Jun 2025}
    \item 1 paper:  \cite{Wasfy2025QARI-OCR:}- \textbf{Sep 2025}
    \item 1 paper:  \cite{Hennara2025Baseer:}- \textbf{Oct 2025}
    \item 1 paper:  \cite{Alhomoud2025Cross-Lingual}- \textbf{Dec 2025}
    \item 1 paper:  \cite{Manna2025Label-Free}- \textbf{Jan 2026}
    \item 3 papers:  \cite{Elkousy2026Domain-Specific} \cite{Taghadouini2026LightOnOCR:} \cite{Malik2026synthocr-gen:}- \textbf{Feb 2026}
    \item 1 paper:  \cite{Faraz2026Designing}- \textbf{Apr 2026}
    \item 2 papers:  \cite{Kargaran2026GlotOCR} \cite{Momayiz2026AtlasOCR:}- \textbf{May 2026}
    \item 4 papers:  \cite{Heakl2026DocAtlas:} \cite{Xu2026Multilingual} \cite{Al-Azzawi2026Cross-Lingual} \cite{Al-Azzawi2026Understanding}- \textbf{Jun 2026}
    \item 1 paper:  \cite{Singh2026Can}- \textbf{Aug 2026}
    \item 1 paper:  \cite{Achyuth2026What}\textbf{Figure 5:} Timeline of low-resource OCR progress; larger markers indicate more citations.
    \end{itemize}
    
    
    \end{itemize}
    
    
    \subsubsection*{3.4 RL, Hard Cases, and System Comparisons}
    
    
    Direct evidence for hard-example mining and RL with verifiable rewards is still limited in low-resource Indic and Arabic OCR, but recent compact OCR VLMs use reward-based post-training to target persistent failure modes that supervised fine-tuning misses  \cite{Taghadouini2026LightOnOCR:}. LightOnOCR-2-1B applies RLVR to repetition loops, formatting failures, and layout-sensitive consistency constraints, then combines gains with checkpoint averaging and task arithmetic  \cite{Taghadouini2026LightOnOCR:}. More broadly, OCR-aware supervised post-training with large synthetic multilingual data reduced hallucination from 18.3\% to 5.5\% and improved OCR completeness from 71.3 to 84.6, suggesting that “hard” or failure-prone cases can be specifically optimized even when CER is not always reported as the main endpoint  \cite{Xu2026Multilingual}.
    
    
    Compared with state-of-the-art multilingual OCR systems from 2024 to 2026, compact adapted models are competitive within their target scripts but not universally dominant. QARI-OCR and Baseer set strong Arabic open benchmarks  \cite{Wasfy2025QARI-OCR:} \cite{Hennara2025Baseer:}. Chitrapathak-2 achieves state-of-the-art or near-state-of-the-art on multilingual Indic benchmarks with better speed than a generic end-to-end VLM  \cite{Faraz2026Designing}. Yet frontier proprietary models still lead many real-scan multilingual evaluations: Gemini was best on Pashto and among the strongest on Devanagari, while Arabic KITAB-Bench found modern VLMs outperforming traditional OCR approaches on average  \cite{Haq2026PsOCR:} \cite{Singh2026Can} \cite{Heakl2025KITAB-Bench:}. At the same time, broad Unicode benchmarks show that all current systems still degrade sharply beyond well-covered scripts, with Arabic and Devanagari already showing substantial drops and most low-resource scripts falling below 10\% acceptance  \cite{Kargaran2026GlotOCR}.
    
    
    \subsubsection*{Results Timeline}
    
    
    Low-resource OCR research has moved from script-specific CRNN systems and synthetic augmentation toward OCR-specialized VLM adaptation, then to compact end-to-end OCR models with post-training for robustness and hallucination control.
    
    
    \subsubsection*{Top Contributors}
    
    
    Based on this corpus, Arabic OCR and multilingual OCR benchmarking are driven by a small recurring set of authors and venues, while arXiv dominates rapid model-release dissemination.
    
    
    \mbox{}\\
    \begin{adjustbox}{max width=\textwidth}
    \begin{tabular}{p{5.3cm} | p{5.3cm} | p{5.3cm}}
    \hline
    Type & Name & Papers \\
    \hline
    Author & E. Barney & [14eb77633a525575b65f3a32eb05fda1][e07b9e846a15579cbbec739c768b0681][f0cf7945d2845ad686a60e7cde30d8f2][9b66763b8a7a54dc81aafb86d98ab450] \\
    Author & M. Liwicki & [14eb77633a525575b65f3a32eb05fda1][f0cf7945d2845ad686a60e7cde30d8f2][9b66763b8a7a54dc81aafb86d98ab450] \\
    Author & C. V. Jawahar & [f8ff9737fe085c22897d5f12c9d3ecc2][4796181bda85500f9c4233b300da31b8][ba8274c8da37560eaf5a043a6793f26a][09511602685e5cbea1208f11266cb68a][a285fb5c1eb054e59e427319f19a8f77] \\
    Journal & \textit{ArXiv} & [012b144be08e5bc18413cd1353232eed][7fe0a5f7b529561cb1dbcdf0e1df09f5][0cba84a76ac75b50b391c37a55d73fad][af798d7c4c845a5c9897f52591645c2a][e00a8bbbbd7f56cab65468f2cdc8aa67][e4d1a82f1ae55540bae517c1a9cafc28][530156972af95e4eb714087b69d33bbd][a453cb9dfd9057858f44c07c7e92c994] \\
    Journal & \textit{IEEE Access} & [a285fb5c1eb054e59e427319f19a8f77][e06176bfdb235120910d341bf4156b7b][2c864665e2c05436b7063a1f46396b38] \\
    Journal & \textit{Pattern Recognition} & [e07b9e846a15579cbbec739c768b0681] \\
    \hline
    \end{tabular}
    \end{adjustbox}
    
    
    \textbf{Figure 6:} Authors and journals most frequent in included papers
    
    
    \subsection*{4. Discussion}
    
    
    The strongest conclusion is that low-resource OCR now benefits substantially from targeted fine-tuning of compact VLMs, but those gains depend more on domain-matched training data and strong OCR initialization than on LoRA or quantization alone  \cite{Faraz2026Designing} \cite{Elkousy2026Domain-Specific} \cite{Hennara2025Baseer:}. Arabic evidence is stronger than Indic evidence because several recent papers report explicit CER or WER improvements after adaptation on shared benchmarks, whereas many Indic papers report WRR, ANLS, or benchmark-relative rank instead of CER  \cite{Wasfy2025QARI-OCR:} \cite{Hennara2025Baseer:} \cite{Faraz2026Designing} \cite{Achyuth2026What}. Real-world robustness remains the field’s main weakness: multiple corpora show that degradation, handwriting variability, diacritics, layout complexity, and label noise still move CER by several points or cause outright failure  \cite{Al-Azzawi2026A} \cite{Al-Azzawi2026Performance} \cite{Farazi2026When}.
    
    
    The labelled-data question has only a partial answer. Cross-lingual Arabic-script HTR shows useful gains with 100, 500, and 1000 labelled lines, while extreme low-resource Indic work treats fewer than 10,000 real images as a meaningful ceiling and supplements them with under 250,000 synthetic images  \cite{Al-Azzawi2026Cross-Lingual} \cite{Achyuth2026What}. Few-shot studies suggest five labelled samples per class can work for character recognition, but full OCR pipelines still appear to need at least hundreds of lines or thousands of word images unless a very strong pretrained OCR model is available  \cite{Chandrashekar2025Applying} \cite{Evani2025SemiHastakshar:}. Annotation quality also matters enough that cleaning labels alone can improve CER by up to 1.8 points in Arabic-script HTR  \cite{Al-Azzawi2026A}.
    
    
    \mbox{}\\
    \begin{adjustbox}{max width=\textwidth}
    \begin{tabular}{p{4.0cm} | p{4.0cm} | p{4.0cm} | p{4.0cm}}
    \hline
    Claim & Evidence Strength & Reasoning & Papers \\
    \hline
    Domain-adapted small or mid-size OCR VLMs reduce Arabic OCR error substantially. & Evidence strength: Strong (8/10) & Multiple 2025–2026 adaptation papers report strong benchmark gains and new open baselines. & \cite{Wasfy2025QARI-OCR:} \cite{Elkousy2026Domain-Specific} \cite{Hennara2025Baseer:} \cite{Momayiz2026AtlasOCR:} \\
    Synthetic data is necessary but insufficient without realism or some real scans. & Evidence strength: Strong (8/10) & Many studies show synthetic pretraining helps, but mixed or realistic data transfers better to real documents. & \cite{Chung2025Finetuning} \cite{Singh2026Can} \cite{Alhomoud2025Cross-Lingual} \cite{Gunda2026CURIO:} \\
    4-bit LoRA or QLoRA usually preserves OCR adaptation quality. & Evidence strength: Moderate (6/10) & Several direct reports show good outcomes, but few controlled 4-bit versus 8-bit OCR ablations exist. & \cite{Wasfy2025QARI-OCR:} \cite{Elkousy2026Domain-Specific} \cite{Momayiz2026AtlasOCR:} \\
    RLVR or GRPO helps compact OCR VLMs reduce hard failures. & Evidence strength: Moderate (4/10) & Evidence is recent and mostly from general OCR benchmarks, not many low-resource Indic or Arabic CER studies. & \cite{Taghadouini2026LightOnOCR:} \cite{Xu2026Multilingual} \\
    Hard-example mining specifically lowers CER in Indic or Arabic OCR. & Evidence strength: Weak (2/10) & The idea is plausible, but direct script-specific OCR evidence remains sparse in this corpus. & \cite{Taghadouini2026LightOnOCR:} \cite{Li2026HunyuanOCR-1.5:} \\
    \hline
    \end{tabular}
    \end{adjustbox}
    
    
    \textbf{Figure 7:} Key claims and supporting evidence in these papers
    
    
    \subsection*{5. Conclusion}
    
    
    Small OCR vision-language models do improve CER for low-resource Indic and Arabic-script tasks when adapted with LoRA or QLoRA, especially if the base model is already OCR-specialized and training mixes realistic synthetic data with at least some real scans. The strongest current evidence is in Arabic, where compact adapted Qwen-based systems reach state-of-the-art open results, while Indic gains are real but more benchmark-dependent and less consistently reported in CER  \cite{Wasfy2025QARI-OCR:} \cite{Hennara2025Baseer:} \cite{Faraz2026Designing} \cite{Singh2026Can}.
    
    
    \subsubsection*{Research Gaps}
    
    
    The literature covers Arabic adaptation, synthetic dataset generation, and multilingual benchmarking reasonably well, but it is thinner on controlled ablations that isolate quantization, hard-example selection, and reinforcement learning effects for specific low-resource scripts.
    
    
    \mbox{}\\
    \begin{adjustbox}{max width=\textwidth}
    \begin{tabular}{p{2.7cm} | p{2.7cm} | p{2.7cm} | p{2.7cm} | p{2.7cm} | p{2.7cm}}
    \hline
     & Arabic & Indic & Mixed Scripts & CER Reported & Data Scaling \\
    \hline
    LoRA Adaptation & GAP & GAP & GAP & GAP & GAP \\
    Synthetic Corpora & GAP & GAP & GAP & GAP & GAP \\
    Real-Scan Benchmarks & GAP & GAP & GAP & GAP & GAP \\
    RL Post-Training & GAP & GAP & GAP & GAP & GAP \\
    Label Cleaning & GAP & GAP & GAP & GAP & GAP \\
    \hline
    \end{tabular}
    \end{adjustbox}
    
    
    Arabic has the clearest adaptation evidence, but even there controlled comparisons between 4-bit, 8-bit, and full fine-tuning are rare  \cite{Elkousy2026Domain-Specific} \cite{Momayiz2026AtlasOCR:}. Indic OCR has more studies on synthetic corpora and multilingual transfer than on compact VLM ablations, and real-scan evaluation remains patchy outside a few stress tests and deployment papers  \cite{Achyuth2026What} \cite{Singh2026Can} \cite{Faraz2026Designing}. Reinforcement learning is the biggest open gap: compact OCR VLMs now use RLVR, but almost none of the low-resource Indic or Arabic papers quantify its CER effect in a controlled script-specific way  \cite{Taghadouini2026LightOnOCR:}.
    
    
    \subsubsection*{Open Research Questions}
    
    
    The most useful next studies would isolate training ingredients rather than introducing new systems without ablations.
    
    
    \mbox{}\\
    \begin{adjustbox}{max width=\textwidth}
    \begin{tabular}{p{8.0cm} | p{8.0cm}}
    \hline
    Question & Why \\
    \hline
    \textbf{How does 4-bit versus 8-bit LoRA affect CER and catastrophic failures on real Arabic and Indic scans?} & Quantized PEFT is widely used, but direct OCR ablations are scarce and deployment decisions depend on this trade-off. \\
    \textbf{What minimum mix of synthetic and real pages is sufficient per script to match larger multilingual OCR systems?} & Current evidence suggests mixed data helps, but the labelled-data threshold remains poorly quantified across scripts. \\
    \textbf{Do RLVR or GRPO reduce CER on low-resource scripts, or mainly reduce hallucination and formatting failures?} & Recent OCR VLMs use reward-based training, yet its script-specific transcription benefit is still unresolved. \\
    \hline
    \end{tabular}
    \end{adjustbox}
    
    
    \textbf{Figure 8:} Open research questions for low-resource OCR adaptation
    
    
    The biggest shift is that compact adapted OCR VLMs are now genuinely competitive, and the biggest open question is how much real script-specific data they still need to stay robust on real scanned pages.
    
    % These search results were found and analyzed using Consensus, an AI-powered search engine for research. Try it at https://consensus.app. © 2026 Consensus NLP, Inc. Personal, non-commercial use only; redistribution requires copyright holders’ consent.
    
    \bibliographystyle{plain}
    \bibliography{refs}
    \end{document}